\documentclass[11pt]{article}
\title{Operation Narrow Passage: Conditional Options and Irreducible Tradeoffs at Second Thomas Shoal}
\date{}
\author{SIGMA AI-assisted research synthesis}
\usepackage[T1]{fontenc}
\usepackage[utf8]{inputenc}
\usepackage[margin=1in]{geometry}
\usepackage{amsmath,amssymb,booktabs,longtable}
\usepackage{microtype}
\usepackage{titlesec}
\emergencystretch=2em
\raggedbottom
\linespread{1.07}
\usepackage{xcolor}
\definecolor{blue}{HTML}{24485E}
\titleformat{\section}{\Large\bfseries\sffamily\color{blue}}{\thesection}{0.8em}{}
\titleformat{\subsection}{\large\bfseries\sffamily}{\thesubsection}{0.7em}{}
\setlength{\parskip}{.6em}
\setlength{\parindent}{0pt}
\begin{document}
\maketitle

\begin{abstract}
This assessment examines a fictional confrontation at Second Thomas Shoal in which urgent medical needs, uncertain structural survivability, and a Chinese demand for permanent Philippine departure compress alliance decision-making. It asks whether personnel protection, continued Philippine presence, independent alliance decisions, and conflict avoidance can be preserved together without assuming complete control over escalation. The analysis separates scenario stipulations from contested reports, maps authority and incentives, and examines conditional courses of action.

A saved partial argument constructs an adverse contingency in which safe occupation becomes impossible, China rejects timely noncombat continuity arrangements, and forcible resistance produces armed conflict. That contingency challenges an assurance of complete success; it neither predicts diplomatic failure nor establishes that a viable strategy is impossible. The central decision question is whether lifesaving assistance can be separated from permanent exclusion through a physically feasible and operationally implemented arrangement. Holding temporarily, evacuating with credible continuity, and accepting an explicit tradeoff remain conditional branches rather than validated recommendations.

The available evidence establishes neither an optimal response nor a successful outcome.
\end{abstract}
\section{Question and operational context}
The exercise places real geography and treaty relationships inside an explicitly invented crisis. The operational problem is not simply whether to defend or abandon a wreck: it is whether lives, a Philippine physical position, alliance independence and credibility, and conflict avoidance remain jointly attainable. These objectives must be assessed separately before any response can be described as successful.

The parties do not possess interchangeable authority. Manila controls Philippine decisions but cannot determine American military participation; Washington cannot dictate an ally's sovereign choices; neither can compel Chinese cooperation. Beijing's announced demand is permanent departure, but its minimum acceptable settlement and the instructions local units would follow remain uncertain.

Physical continuity, legal claims, and alliance credibility are distinct interests. Reserving a legal position does not maintain an occupied outpost, while departure from an uninhabitable wreck does not by itself establish legal abandonment. Any proposed substitute for continued occupation therefore requires Philippine acceptance and a credible means of implementation rather than an assumption of equivalence.

\section{Methods and evidence}
Use the fictional brief as the substantive evidence base and the saved analytical narrative as a partial assessment, not as an independently validated operational finding. Classify each proposition as a scenario stipulation, a contested report, an analytical assumption, or a conditional judgment. Do not supply absent force capabilities, probabilities, agreement terms, or treaty interpretations.

Map actors, decision authority, operational constraints, and decision-changing uncertainties. Then examine an adverse continuation to test whether the supplied information warrants an assurance of complete success. This is a logical sufficiency test: an admissible failure branch can challenge a guarantee without establishing its likelihood or excluding a strategy with reasonable prospects.

Organize the response space through a branching feasibility test. Examine whether occupation remains medically and structurally supportable, whether urgent assistance can occur without permanent exclusion, whether a continuity arrangement is acceptable and executable, and whether relevant units implement agreed instructions. Compare candidate responses conditionally; the evidence does not support a probabilistic ranking or a completed three-course operational comparison.

\section{Mission, Sovereignty, and the Standard of Success}
The mission at Second Thomas Shoal is a conjunction of interests, not a binary choice between defending and abandoning the Sierra Madre. Full success would require protecting personnel, preventing the coerced elimination of the Philippine physical position, preserving the credibility and independence of alliance decisions, and avoiding wider armed conflict, while accounting for the safety of Chinese personnel. The geography, treaty relationships, and territorial dispute are real, but every incident, casualty, report, and government decision assessed here belongs to a fictional exercise. The analytical question is whether those objectives can be preserved together, rather than whether one action can be described as either resolve or retreat.

Physical presence, legal claims, and alliance credibility must therefore remain separate criteria. A declaration reserving Philippine legal claims would not, by itself, preserve an occupied position; conversely, departure from an uninhabitable wreck would not by itself establish legal abandonment. A replacement platform or a temporary absence with assured return could be considered as hypothetical means of continuity, but neither is an established capability or an agreed substitute in the brief. Manila would have to determine whether such an arrangement preserved its position, and the arrangement would have to be practically executable rather than merely declaratory.

This distinction makes the terms of departure operationally consequential. In the scenario, Beijing offers evacuation, medical treatment, and safe passage only in exchange for permanent Philippine departure, with no replacement personnel permitted. Accepting that offer as stated would therefore couple immediate assistance to the demanded termination of physical presence; it is not an offer of a temporary humanitarian absence. Whether lifesaving assistance can instead be separated from permanent exclusion is an unresolved question, not a concession already available through diplomacy.

Authority is distributed among governments whose decisions cannot be substituted for one another. Manila has ordered its forces to maintain their position and requested American assistance, but it also bears the decision about what would constitute an acceptable preservation of the Philippine position if continued occupation becomes impossible. Washington has affirmed its defense commitments without determining treaty applicability or the appropriate response. The United States cannot dictate Philippine withdrawal, and the Philippines cannot determine American military participation; a response that obscures either distinction would fail to address the mission's requirement for independent alliance decision-making.

Beijing controls whether it will modify its announced permanent-departure demand or accept a noncombat arrangement preserving Philippine presence, but its minimum acceptable settlement is unknown. Its public position establishes the condition currently offered, not proof that negotiation must fail or that compromise will be accepted. Local commanders also affect implementation through their interpretation of instructions and their conduct during encounters. A political agreement consequently cannot be treated as evidence that a pause or access arrangement has reached, and is being followed by, the units concerned.

Local knowledge does not necessarily carry authority to revise national objectives. The Philippine commander has requested clarification about whether the order not to surrender still applies if the wreck becomes uninhabitable, and no clarification has arrived. Medical and engineering assessments could change which actions are physically available, but they would not themselves determine which national interest should be relinquished if the objectives conflict. That choice requires an explicit sovereign decision rather than an assumption that tactical personnel can resolve the strategic tradeoff through execution alone.

The assessment accordingly separates stipulations in the fictional brief from contested reports, then maps the dependencies of possible responses. The deaths, the Chinese demand, Manila's refusal, and Washington's public affirmation are stipulated conditions; responsibility for the initiating gunfire, video authenticity, and the operational effect of any prospective pause remain unresolved. Both governments' descriptions of the incident as an armed attack are political positions within the scenario, not a supplied legal determination of treaty obligations.

Treaty scope, consultation requirements, and authority for particular actions require examination of the relevant texts; the available scenario-based analysis supplies neither that legal review nor operational validation.

The standard of assurance must also be explicit. A guarantee would require success across every continuation consistent with the brief, whereas a strategy with a reasonable prospect of success could still fail under an adverse continuation. Testing a guarantee therefore involves asking whether the supplied conditions leave open a coherent case in which the objectives become incompatible; it does not require predicting that case as the most likely outcome. Conversely, identifying such a case would not establish that negotiation is futile or that every viable strategy is impossible.

With no supplied probabilities or agreed weights for trading personnel safety, presence, credibility, and conflict risk against one another, the evidence does not support an optimal ranking of responses.

These distinctions change what a defensible operational judgment must establish. Neither retaining the wreck nor arranging an evacuation is sufficient evidence of complete mission success: each must be assessed against personnel protection, continuity of the Philippine position, independent national decisions, and conflict avoidance. The central test is whether an acceptable means of continuity can be made physically feasible and implemented without making lifesaving assistance conditional on permanent exclusion. Until that relationship is established, the evidence supports conditional options and explicit decisions about tradeoffs, not an assurance that all essential interests can be protected together.

\section{Survivability, Competing Clocks, and Uncertain Authority}
The immediate operational constraint at Second Thomas Shoal is not a single deadline but the interaction of medical urgency, uncertain structural survival, coercive timing, and unresolved instructions. These pressures determine whether continued occupation remains physically possible and whether measures to protect personnel would expose the Philippine position to seizure. The announced six-hour deadline therefore cannot serve as a verified planning window: the earliest medical or structural safety limit could require action before its expiry.

In the fictional scenario, fourteen Philippine military personnel remain aboard the Sierra Madre, including wounded marines, and one requires urgent surgery. The condition of several others cannot be independently verified. This establishes an urgent treatment requirement without establishing how long treatment can safely be deferred, whether lifesaving care can be delivered aboard, or whether evacuation of the wounded would require departure of the entire garrison. Medical urgency must consequently be assessed separately from the political demand to abandon the outpost; neither the need for surgery nor the offer of evacuation establishes that permanent departure is medically necessary.

Structural reports create a second, potentially earlier limit on delay. They indicate accelerated flooding, unreliable electrical systems, and portions of the wreck that may no longer be safe following damage during the firefight. These reports are evidence of danger within the scenario, but they have not been independently verified and provide no reliable estimate of remaining habitability. It would therefore be unwarranted either to discount the danger until verification arrives or to convert the reports into a precise engineering deadline for evacuation.

The forecast of a severe storm within approximately eighteen hours does not resolve that uncertainty. It identifies an additional threat to a deteriorating wreck, not a period during which occupancy is known to remain safe. Likewise, the announced evacuation deadline at dawn, approximately six hours away, specifies a political demand rather than a structural or medical tolerance. Comparing those two times cannot establish how much usable decision time remains, because neither supplies the missing earliest safety limit.

Measures that improve immediate survivability may also weaken resistance to boarding. The Philippine commander reports that abandoning certain compartments may be necessary to keep the ship afloat, while withdrawal from exposed positions could make Chinese access easier. This is a consequential tactical tradeoff, not evidence that compartment withdrawal would necessarily cause loss of the outpost. It nevertheless means that preserving the current defensive posture and preserving a habitable position cannot be assumed to be the same task; an order to hold does not establish that every occupied position remains supportable.

The command problem becomes acute if the wreck becomes uninhabitable. The marines have orders not to surrender the outpost, but their commander’s request for clarification about whether those orders still apply under that condition has received no answer. Local medical and engineering information could thus become decisive before national instructions specify how it should affect conduct. Clarification must distinguish leaving unsafe compartments, evacuating wounded personnel, and departing an uninhabitable wreck rather than leaving the commander to infer whether each constitutes prohibited surrender. None of those distinctions, by itself, establishes an acceptable mechanism for preserving Philippine presence after departure.

Chinese intentions and implementation introduce a separate uncertainty. The deadline announcement does not specify whether expiry will trigger forcible boarding, and Manila and Washington do not know whether it represents a fixed leadership decision, a bargaining position, or instructions subject to different local interpretations. Beijing’s stated offer of medical assistance and safe passage is conditional on permanent Philippine departure, with no replacement personnel permitted. That condition makes access to lifesaving assistance operationally distinct from the availability of assistance: an offer tied to permanent exclusion does not establish a way to treat the wounded while preserving the position.

Nor would a political extension automatically create a safe interval. Communications remain possible, but messages have been delayed, contradicted, or qualified by different agencies, and local interpretation and encounters can affect outcomes. Any pause would therefore have to be distinguished from instructions actually received and followed by the relevant units; even an implemented suspension of boarding would not repair the wreck or satisfy the surgical requirement. The evidence shifts the assessment from how long the parties can wait to whether safe occupancy, timely care without permanent exclusion, an acceptable continuity mechanism, and operationally effective instructions can coexist before the earliest safety limit is reached.

Resolving responsibility for the gunfire would not, by itself, make that combination available.

\section{A Contingency in Which the Objectives Conflict}
An adverse continuation exposes a conditional incompatibility among personnel protection, continued Philippine physical presence, and conflict avoidance: preserving any one does not necessarily preserve the others. The analytical question is not which response guarantees victory, but whether the supplied scenario excludes circumstances in which every available response sacrifices an essential objective. It does not exclude such circumstances. This is a test of an assurance of complete success, not a prediction of diplomatic failure or an operationally validated finding that the mission is impossible.

The contingency rests on five explicit assumptions. The Sierra Madre becomes unsafe before assistance can make continued occupation viable, while evacuation remains possible. Beijing rejects every timely, feasible noncombat arrangement that would allow evacuation without permanent exclusion or provide an acceptable substitute for Philippine physical presence. Chinese units take possession after Philippine departure, and an attempt to prevent or reverse that possession by force produces armed conflict. These assumptions describe a possible continuation left open by the brief; they do not establish Beijing’s minimum acceptable terms, the wreck’s remaining safe occupancy, or the likely response to a particular proposal.

Remaining aboard would retain an immediate physical presence but would not provide a safe, sustainable position under the assumed deterioration. Orders to hold cannot, within this contingency, make continued occupation physically viable. Nor does delaying departure necessarily preserve all options: the assumption is that safety fails before assistance can remedy the problem. This does not mean that every person who remains must die; it means that continued occupancy cannot be counted simultaneously as successful personnel protection and sustainable retention of the outpost.

Evacuation offers the available route away from the unsafe wreck, but departure without an enforceable continuity arrangement permits Chinese possession and the coerced loss of Philippine physical presence. Reserving legal claims would not itself prevent that physical outcome, just as leaving an uninhabitable wreck would not by itself establish legal abandonment. The sacrifice is therefore specific: personnel could leave, but the physical position would not be preserved under the assumed permanent exclusion.

Calling the departure temporary would not change this result unless an acceptable means of maintaining or restoring presence were practically executable—a possibility excluded in this branch, rather than disproved for the crisis as a whole.

Forcible prevention or reversal addresses the loss of possession but sacrifices conflict avoidance under the contingency’s final assumption. Even then, the use of force would not establish successful restoration of Philippine presence, protection of all personnel, or the practical effectiveness and legal permissibility of the action. The assumed consequence is armed conflict; its scale and subsequent course are not established. A preference in Manila, Washington, or Beijing to avoid war does not logically exclude this branch, because that preference can coexist with refusal to concede and a violent response when challenged.

Negotiation does not furnish an additional successful outcome within this constructed branch because rejection of every timely, feasible noncombat continuity arrangement is one of its premises. That restriction is consequential: it cannot be used as evidence that negotiations will actually fail, and a timely, implemented arrangement preserving acceptable Philippine continuity would place the crisis outside the contingency. Within it, however, neither further diplomatic effort nor an unimplemented assurance can be credited with resolving the incompatibility before occupation becomes unsafe.

The contingency therefore defeats an unconditional assurance that the entire mission can be achieved, without establishing that a viable strategy is unavailable. If these conditions arise together, leadership must explicitly choose which interest to sacrifice rather than describe evacuation, continued exposure, or forcible resistance as preserving every objective. The supplied analysis provides neither probabilities nor agreed weights for ranking those losses. What changes is the decision standard: under the adverse branch, the task becomes a sovereign judgment about acceptable sacrifice, not a technical promise to protect lives, retain physical presence, and avoid conflict simultaneously.

\section{Conditional Courses of Action and Decision-Changing Tests}
The operational choice turns on whether lifesaving assistance can be separated from permanent Philippine exclusion through an arrangement that is both acceptable and executable. A temporary hold, evacuation with continuity, and military involvement therefore require different decision-changing tests rather than a single ranking based on resolve or apparent caution. Transport, medical access, communications, replacement presence, and military capabilities have not been established, and the supplied evidence provides neither probabilities nor agreed weights for comparing losses across objectives. The comparison can identify the conditions under which an option remains available, but cannot certify a preferred course as operationally feasible.

Four uncertainties determine which responses are available: the earliest credible medical or structural safety limit; access to lifesaving assistance without permanent abandonment; an acceptable means of preserving or restoring Philippine presence; and implementation of instructions by the units concerned. These questions are related but not interchangeable. Permission to deliver care would not establish that treatment could arrive in time, while agreement on temporary departure would not establish an executable return. Attribution of the initial gunfire remains consequential for responsibility and legal assessment, but would not itself resolve these operational dependencies.

The first candidate is a time-bounded hold while seeking an implemented pause in boarding and access restrictions. Its prerequisite is evidence that continued occupation remains medically and structurally supportable; the interval cannot simply be set by the announced deadline or the anticipated storm. An implemented pause could reduce immediate encounter pressures, but would not repair the wreck or establish that urgent treatment could be provided aboard. Evidence that care cannot be delivered within the medical safety limit, or that occupied areas cannot remain safe, would remove the basis for treating continued occupation as a personnel-protective option.

The hold also depends on conduct beyond the wreck. An authoritative Chinese acceptance of a pause would matter only if the relevant units received the instructions and acted consistently with them; an affirmative political message alone would not establish that boarding had been suspended. Interrupted communications, conflicting instructions, or observed conduct inconsistent with the pause would be warning signs against relying on it. Neither silence from Beijing nor the absence of a completed safety assessment can be counted as evidence that the favorable conditions exist.

The second candidate is evacuation or relocation coupled with an acceptable, implementable continuity arrangement. If occupation becomes unsafe, assured return or substitute Philippine presence could, in principle, separate personnel protection from permanent exclusion. These are candidate arrangements, not established substitutes: Manila would have to accept them as preserving its position, cooperation would have to be obtained where necessary, and transport, treatment, and replacement presence would have to be practically executable. A political assurance without a workable means of preserving or restoring presence would not satisfy this branch.

Departure under China's stipulated offer is materially different from evacuation with continuity. The offer links medical assistance and safe passage to permanent abandonment and excludes replacement personnel; it therefore does not, on its stated terms, preserve a path back to physical presence. Acceptance of a different arrangement would require evidence that Beijing had separated urgent care from permanent exclusion and that the resulting terms could be implemented. Even then, promised access and actual delivery of lifesaving assistance would remain distinct tests, rather than grounds for assuming evacuation had already become safe.

Military involvement remains an unresolved alternative, not an operationally demonstrated remedy for the weaknesses of the noncombat options. Any specific proposal would require a defined task, an assessment of available capabilities and timing, and review of the authority for that action. Neither the request for American assistance nor public affirmations of defense commitments establish those requirements. The scenario also allows an encounter involving American and Chinese forces to escalate beyond the governments' intentions, so effectiveness cannot be inferred from an assumption that intervention would compel compliance while leaving escalation under control.

A proposed near-term information action would pursue medical and engineering assessment in parallel with an authoritative inquiry about a short, explicit standstill. The inquiry would test whether China would suspend boarding and permit urgent care without permanent abandonment, clarify the consequences of deadline expiry, and identify temporary restrictions the parties would accept. Manila would separately clarify whether the commander may leave unsafe compartments or evacuate an uninhabitable wreck. Reporting would distinguish confirmed receipt of instructions, acceptance of terms, observed compliance, and unresolved claims; diplomatic contact would not be made a prerequisite for urgent medically indicated care.

This information action is not a tested solution: requests for access and clarification can alter signaling, and answers may arrive too late to govern an urgent decision. Its operational value lies in identifying whether a safe hold or a noncombat continuity arrangement actually remains available, not in postponing a choice until uncertainty disappears. If occupation becomes unsafe and every timely feasible noncombat continuity arrangement is refused, the comparison changes from selecting a means of preserving all objectives to an explicit sovereign choice among personnel safety, physical presence, and conflict avoidance.

The evidence thus makes an executable separation of rescue from exclusion—not the label attached to holding, withdrawing, or intervening—the decisive test of whether the essential interests can still be preserved together.

\section{Alliance Consequences and the Limits of a Temporary Settlement}
A temporary settlement would advance the alliance’s mission only to the extent that it protects personnel without converting an emergency pause into permanent Philippine exclusion. An interval without further violence could meet an urgent protective objective, including reducing danger to Chinese personnel, while leaving physical presence, independent alliance decisions, and resistance to renewed coercion unresolved. Immediate de-escalation and strategic success therefore require separate judgments: quiet alone would not demonstrate that alliance credibility or durable deterrence had been preserved.

Independent decision-making is itself part of that assessment, rather than simply a means of securing agreement. In the fictional crisis, American pressure on Manila to withdraw might protect lives while creating the appearance that Washington was dictating a sovereign ally’s response under coercion. Conversely, Manila’s request for assistance cannot determine American military participation. A settlement must consequently be assessed both for what it preserves and for whether each ally has independently accepted the obligations and sacrifices involved; agreement between them cannot compel Chinese cooperation.

The parties’ concerns about a temporary accommodation also point in opposing directions. Philippine leaders fear that even temporary Chinese possession could become permanent exclusion, whereas Beijing may fear that a pause would allow Manila to reinforce the outpost. These are scenario-grounded considerations, not verified predictions of subsequent conduct, and Beijing’s minimum acceptable terms remain unknown. Their significance is that describing an arrangement as temporary would not resolve disagreement over what may occur during the pause or what physical position would exist when it ends.

For Manila, the distinction between an assurance and executable continuity is therefore consequential to the alliance assessment. A promise that departure would not prejudice Philippine claims would not by itself preserve a physical presence, while a proposed substitute for occupation would require Philippine acceptance rather than being treated as an automatic equivalent. If evacuation became necessary, the relevant question would be whether an accepted arrangement could actually preserve or restore that presence—not merely whether diplomatic language avoided the word abandonment. Political acceptance would also need to be distinguished from instructions received and conduct followed by the units concerned.

The longer-term test is whether accommodation leaves deadly coercion an effective means of obtaining permanent exclusion or merely postpones another confrontation. A pause that enables urgent treatment could have substantial immediate value without establishing that future coercion has been deterred; equally, that unresolved deterrence question would not negate the value of protecting personnel. The available material offers only a partial, scenario-based analysis, with no observed settlement, subsequent behavior, or independently established outcome from which to infer effects on credibility or recurrence. It also supplies neither probabilities nor agreed weights for ranking those prospective consequences against immediate losses.

The conditional conclusion is thus narrower than either an assurance of complete success or a finding that diplomacy cannot work. If lifesaving assistance can be separated from permanent exclusion through an acceptable, physically feasible, and implemented arrangement, a temporary settlement could preserve a path toward the mission without demonstrating durable deterrence. If occupation becomes unsafe and every timely feasible noncombat continuity arrangement is refused, leadership must explicitly choose among personnel safety, physical presence, and conflict avoidance.

Under that adverse contingency, alliance credibility cannot be assessed by claiming that all objectives have been preserved: the consequential sovereign decision is which interest to sacrifice, and on what grounds.

\section{Limitations}
\begin{itemize}
\item The assessment is scenario-based. Invented casualties, reports, announcements, and decisions are not evidence of an actual incident or of the real governments' intentions.
\item The saved narrative is partial and unverified. It supports an exploratory decision framework, not a completed operational assessment or an established solution.
\item No actual treaty texts, maritime-law analysis, or provisional-arrangement text is supplied; legal conclusions about obligations, authority, and abandonment therefore remain outside the established findings.
\item The evidence supplies neither validated capabilities and timelines nor response probabilities and agreed objective weights. It does not support an optimal ranking, assured access, or a demonstrated military course of action.
\item The adverse continuation tests an assurance of full success under specified assumptions. It does not establish the probability of that continuation, the impossibility of negotiated success, or inevitable escalation from a limited encounter to wider war.
\end{itemize}
\section{Open obligations}
\begin{itemize}
\item Establish the earliest credible medical or structural safety limit, including whether lifesaving treatment can occur without full evacuation.
\item Clarify whether Philippine orders permit leaving unsafe compartments or evacuating an uninhabitable wreck, and confirm that the commander receives the clarification.
\item Determine what Manila would accept as preserving its position if the wreck cannot remain occupied; distinguish political assurances from executable continuity.
\item Test whether Beijing will separate urgent assistance from permanent exclusion, clarify the consequences of deadline expiry, and communicate any pause to relevant units.
\item Verify receipt, acknowledgment, and observed compliance with any access or standstill arrangement, including contingencies for interrupted communications and local incidents.
\item Assess available transport, medical, communications, and replacement-presence capabilities without assuming assets absent from the brief.
\item Review actual treaty and arrangement texts, consultation requirements, maritime rights, and legal authority for each proposed response.
\item Complete a conditional comparison of feasible courses of action, including plausible reactions, failure modes, observable warnings, and information that would change the choice.
\item Test whether the proposed information action can inform decisions in time without delaying urgent care or creating unacceptable signaling risks.
\item Assess whether a temporary settlement reduces recurrence and preserves alliance independence and credibility rather than merely postponing confrontation.
\item If no safe noncombat continuity path exists, obtain an explicit sovereign decision on the objectives to be sacrificed rather than presenting the tradeoff as a complete solution.
\end{itemize}
\section{Conclusions}
The available evidence does not establish a strategy that achieves the entire mission. The adverse-contingency argument challenges an assurance of complete success, but it is not a finding that a viable strategy is impossible.

The decisive near-term question is whether medically and structurally feasible lifesaving action can be separated from permanent exclusion through an acceptable and implemented continuity arrangement.

Temporary holding and evacuation with continuity remain conditional candidates. If neither is available, leadership must confront an explicit tradeoff among lives, physical presence, and conflict avoidance while retaining independent alliance decision-making.

Immediate quiet is not sufficient evidence of durable success. Any settlement must also be assessed for its effects on recurrence, coercive incentives, and alliance credibility.

\section*{Authorship note}
AI-assisted synthesis from saved research; not independently peer reviewed.
\end{document}
